% Provenance (never printed): NSI terminale "Architectures materielles, systemes d'exploitation et reseaux:
% protocoles de routage (identifier, suivant le protocole de routage utilise, la route empruntee par un paquet);
% RIP (distance vector, hop count) et OSPF (link state, cost from bandwidth)". See sources/book2/BRIEFS.md.
\chapter{Routing}\label{ch:g12:routing}

A digger cuts a cable between two routers of a network. The packets that used it do not vanish: within seconds or
minutes, the routers on either side notice the silence, tell their neighbours, and the packets take another way.
Nobody configured that other way by hand. The routers found it themselves, by talking to each other.

This chapter shows the two ways routers do it. In the first, each router knows only distances, and learns them from
its neighbours; in the second, each router learns the whole map and computes its own routes with \omterm{met:g12:graph-algorithms:dijkstra}{Dijkstra's
algorithm}. Both are simulated on one invented network of seven routers, before and after a link fails, and the first
shows a strange failure of its own: \omterm{def:g12:routing:convergence}{counting to infinity}.

\begin{recall}
A router passes each packet to a \emph{\omterm{def:g11:network-protocols:routing-table}{next hop}} found in its \omterm{def:g11:network-protocols:routing-table}{routing table}; the whole route is the chain of these
local choices (chapter \ref{ch:g11:network-protocols}). A \omterm{def:g10:graphs-first-look:directed}{weighted graph} has a number on each edge
(chapter \ref{ch:g10:graphs-first-look}). \omterm{met:g12:graph-algorithms:dijkstra}{Dijkstra's algorithm} finds the shortest paths from one vertex when the
weights are not negative (chapter \ref{ch:g12:graph-algorithms}).
\end{recall}

\section{A network as a weighted graph}\label{sec:g12:routing:graph}

The routers of a network are the vertices of a \omterm{def:g10:graphs-first-look:graph}{graph}, and the links between them its edges. Each link has a
\emph{rate}, the number of bits it carries per second. To choose routes, each link also gets a number, its cost.

\begin{definition}[Hop count, route cost]\label{def:g12:routing:metric}
The \emph{hop count}\index{hop count} of a route is its number of links. When each link has a cost, the \emph{route
cost}\index{route cost} is the sum of the costs of its links. A routing method chooses, for each destination, a route
whose hop count, or whose cost, is smallest.
\end{definition}

The \omterm{def:g12:routing:metric}{hop count} treats every link alike. A cost can prefer fast links: this chapter gives each link the cost
$10\,000$ divided by its rate in megabits per second, so a link of $10\,000$ Mbit/s costs $1$, one of $1000$ costs
$10$, one of $10$ costs $1000$. The protocol that uses such costs leaves each link's cost to the network's
administrator; dividing a reference rate by the link's rate is one way to set them.
% ledger: g12:ospf-2328

\omcode{code/grade-12/11-routing/listings/network.py}{1}{33}{The network of the chapter, and the two ways of costing
its links.}\label{lst:g12:routing:network}

\begin{omfigure}
\begin{tikzpicture}[font=\footnotesize, r/.style={circle, draw, thick, fill=omDef!12, minimum size=8mm},
    l/.style={thick}, lab/.style={fill=white, inner sep=1.5pt, align=center, font=\scriptsize}]
  \node[r] (A) at (0,0) {A};
  \node[r] (B) at (3,1.5) {B};
  \node[r] (C) at (3,-1.5) {C};
  \node[r] (D) at (6.5,1.5) {D};
  \node[r] (E) at (6.5,-1.5) {E};
  \node[r] (F) at (10,1.5) {F};
  \node[r] (G) at (10,-1.5) {G};
  \foreach \u/\v/\rate/\cost/\w in {A/B/1000/10/very thick, A/C/100/100/thick, B/C/1000/10/very thick,
      B/D/100/100/thick, C/E/10/1000/thin, D/E/1000/10/very thick, D/F/1000/10/very thick,
      E/G/100/100/thick, F/G/1000/10/very thick}
    \draw[\w] (\u) -- node[lab] {\rate\\[-1pt] cost \cost} (\v);
\end{tikzpicture}

{\small The invented network: seven routers, nine links, each labelled with its rate in Mbit/s and its cost. Thicker
lines are faster links. By \omterm{def:g12:routing:metric}{hop count}, the best route from A to E goes through C, over the slowest link of all.}
\end{omfigure}

\omcode{code/grade-12/11-routing/listings/ospf_costs.py}{1}{14}{The cost of each link, and the route from A to E by
each measure.}\label{lst:g12:routing:costs}

\omoutput{code/grade-12/11-routing/outputs/ospf_costs.txt}{Counting hops sends the traffic over a 10 Mbit/s link;
costs from rates send it one link further, but ten times faster at its
slowest.}\label{out:g12:routing:costs}

\begin{definition}[Routing protocol, distance-vector routing, link-state routing]\label{def:g12:routing:routing-protocol}
A \emph{routing protocol}\index{routing protocol} is a set of rules by which routers exchange messages so that each
can fill its own \omterm{def:g11:network-protocols:routing-table}{routing table}. In \emph{distance-vector routing}\index{distance-vector routing}, each router sends
its neighbours its table of distances, and computes its own from theirs. In \emph{link-state
routing}\index{link-state routing}, each router sends every other router the list of its own links and their costs;
each router then knows the whole \omterm{def:g10:graphs-first-look:graph}{graph}, and computes its routes from it.
\end{definition}

\section{Distance vectors}\label{sec:g12:routing:dv}

A router at first knows only itself. Each \emph{round}, every router receives its neighbours' tables. If neighbour
$n$, reached at cost $c$, says it reaches destination $d$ at distance $x$, then the router can reach $d$ at $c + x$
through $n$. For each destination it keeps the smallest of these sums, and $n$ becomes the \omterm{def:g11:network-protocols:routing-table}{next hop}. This is the
update of the Bellman--Ford algorithm, run by every router at once on its own data.

\omcode{code/grade-12/11-routing/listings/dv_rounds.py}{1}{37}{Distance vectors, round after round, from tables that
know nothing.}\label{lst:g12:routing:dv}

\omoutput{code/grade-12/11-routing/outputs/dv_rounds.txt}{A's distances, round by round. News travels one link per
round; a first answer, such as $1100$ for E, can be bettered later.}\label{out:g12:routing:dv}

After round $k$, a router knows the shortest route to each destination among the routes of at most $k$ links
(exercise \ref{exo:g12:routing:12}). The routes here have at most four links, so four rounds suffice; a fifth changes
nothing, and the tables are stable.

\begin{definition}[Convergence, counting to infinity]\label{def:g12:routing:convergence}
A \omterm{def:g12:routing:routing-protocol}{routing protocol} has \emph{converged}\index{convergence} when no table changes any more and every route is a
shortest one. \emph{Counting to infinity}\index{counting to infinity} is a slow convergence of \omterm{def:g12:routing:routing-protocol}{distance-vector
routing} after a failure: routers learn from each other distances that pass through themselves, and raise them a
little each round until they reach the value that means ``unreachable''.
\end{definition}

\section{Bad news travels slowly}\label{sec:g12:routing:infinity}

Take three routers on a line, A -- B -- C, counting hops. The stable tables say that B reaches C in $1$ hop and A in
$2$, through B. Then the link B -- C fails. B no longer hears from C; but A still says ``I reach C in $2$''. B cannot
know that A's route goes through B itself, so it believes it: $3$, through A. A then hears $3$ from B, its \omterm{def:g11:network-protocols:routing-table}{next hop},
and must update to $4$. The two count up together.

To stop the count, the hop-count protocol sets a ceiling: a distance of $16$ means unreachable, and routes may have at
most $15$ hops.
% ledger: g12:rip-2453-metric

\omcode{code/grade-12/11-routing/listings/count_to_infinity.py}{1}{39}{The line A -- B -- C, before and after the link
B -- C fails.}\label{lst:g12:routing:infinity}

\omoutput{code/grade-12/11-routing/outputs/count_to_infinity.txt}{Each round, one of the two raises its distance by
two, using the other's old value.}\label{out:g12:routing:infinity}

\begin{omfigure}
\begin{tikzpicture}[font=\footnotesize, r/.style={circle, draw, thick, fill=omDef!12, minimum size=7mm}]
  \node[r] (A) at (0,0) {A}; \node[r] (B) at (2,0) {B}; \node[r] (C) at (4,0) {C};
  \draw[thick] (A) -- (B);
  \draw[thick, dashed, omThm] (B) -- node[omThm] {$\times$} (C);
  \draw[->, omMeth, thick] (A) to[bend left=35] node[above, font=\scriptsize] {``C at 2''} (B);
  \draw[->, omThm, thick] (B) to[bend left=35] node[below, font=\scriptsize] {``C at 3''} (A);
  \node[anchor=west] at (5.2,0) {
    \begin{tabular}{l*{9}{c}}
      round & 0 & 1 & 2 & 3 & 4 & \dots & 13 & 14 & 15 \\ \hline
      A & 2 & 2 & 4 & 4 & 6 & \dots & 14 & 16 & 16 \\
      B & 1 & 3 & 3 & 5 & 5 & \dots & 15 & 15 & 16
    \end{tabular}};
\end{tikzpicture}

{\small \omterm{def:g12:routing:convergence}{Counting to infinity}. After the failure each router takes the other's distance, plus one, as its own; the
values climb by two every two rounds until both reach $16$, unreachable, at round $15$.}
\end{omfigure}

The cure for this case is to stop A from telling B a route that goes through B. With \emph{split horizon}, a router
leaves out of the table it sends to a neighbour the routes whose \omterm{def:g11:network-protocols:routing-table}{next hop} is that neighbour. A variant, \emph{poisoned
reverse}, sends those routes with the distance $16$ instead.
% ledger: g12:rip-2453-split

\omcode{code/grade-12/11-routing/listings/split_horizon.py}{1}{3}{The same failure, with split
horizon.}\label{lst:g12:routing:split}

\omoutput{code/grade-12/11-routing/outputs/split_horizon.txt}{B learns nothing about C any more, and marks it
unreachable at once; A follows one round later.}\label{out:g12:routing:split}

Split horizon prevents a loop of two routers. A loop of three or more can still count to infinity; the ceiling of
$16$ is what finally ends it.

\section{Link state}\label{sec:g12:routing:link-state}

In \omterm{def:g12:routing:routing-protocol}{link-state routing}, each router sends its own list of links, with their costs, to its neighbours. A router that
receives a list it has not seen yet keeps a copy and passes it on to all its neighbours except the one it came from;
a list it has already seen is dropped. This is \emph{flooding}: every list reaches every router. Each router then
holds the same map, and runs \omterm{met:g12:graph-algorithms:dijkstra}{Dijkstra's algorithm} from itself; the first router on each shortest path is the \omterm{def:g11:network-protocols:routing-table}{next
hop}.
% ledger: g12:ospf-2328

\omcode{code/grade-12/11-routing/listings/link_state.py}{1}{20}{Flooding the link lists: each router passes on
the lists it has not seen yet.}\label{lst:g12:routing:flood}

\omcode{code/grade-12/11-routing/listings/link_state.py}{23}{44}{Dijkstra from each router, keeping the first router
of each path.}\label{lst:g12:routing:link-state}

\omoutput{code/grade-12/11-routing/outputs/link_state.txt}{The same distances as the distance vectors found, computed
by each router on its own.}\label{out:g12:routing:link-state}

Flooding one list costs one message on each link in each direction, except the links by which a router first heard
the list: here $2 \times 9 - 6 = 12$ messages per list, $84$ for the seven. That is more traffic than one round of
distance vectors, but it is sent only when a link changes.

\begin{omfigure}
\begin{tikzpicture}[font=\footnotesize, r/.style={circle, draw, thick, fill=omDef!12, minimum size=8mm},
    lab/.style={fill=white, inner sep=1pt, font=\scriptsize}]
  \node[r, fill=omMeth!25] (A) at (0,0) {A};
  \node[r] (B) at (3,1.5) {B};
  \node[r] (C) at (3,-1.5) {C};
  \node[r] (D) at (6.5,1.5) {D};
  \node[r] (E) at (6.5,-1.5) {E};
  \node[r] (F) at (10,1.5) {F};
  \node[r] (G) at (10,-1.5) {G};
  \foreach \u/\v/\c in {A/C/100, C/E/1000, E/G/100}
    \draw[black!35, dashed] (\u) -- node[lab, text=black!50] {\c} (\v);
  \foreach \u/\v/\c in {A/B/10, B/C/10, B/D/100, D/E/10, D/F/10, F/G/10}
    \draw[->, ultra thick, omMeth] (\u) -- node[lab, text=black] {\c} (\v);
  \foreach \v/\d in {B/10, D/110, F/120} \node[above=4mm, font=\scriptsize] at (\v) {\d};
  \foreach \v/\d in {C/20, E/120, G/130} \node[below=4mm, font=\scriptsize] at (\v) {\d};
\end{tikzpicture}

{\small A's shortest-path tree, with the cost of each link and A's distance to each router. Every route from A leaves
through B, which is why B is the \omterm{def:g11:network-protocols:routing-table}{next hop} for every destination in A's table. The dashed links are used by no
shortest route from A.}
\end{omfigure}

\begin{omfigure}
\small
\begin{tabular}{lrc@{\hspace{3em}}lrc}
  \multicolumn{3}{c}{\textbf{table of A}} & \multicolumn{3}{c}{\textbf{table of E}} \\
  destination & cost & \omterm{def:g11:network-protocols:routing-table}{next hop} & destination & cost & \omterm{def:g11:network-protocols:routing-table}{next hop} \\ \hline
  B & 10 & B & A & 120 & D \\
  C & 20 & B & B & 110 & D \\
  D & 110 & B & C & 120 & D \\
  E & 120 & B & D & 10 & D \\
  F & 120 & B & F & 20 & D \\
  G & 130 & B & G & 30 & D
\end{tabular}

\medskip
{\small Two \omterm{def:g11:network-protocols:routing-table}{routing tables} side by side, from the output above. A packet from A to G goes to B; B's table sends it on
to D, D's to F, F's to G. E sends even its packets for C through D, avoiding its own slow link to C.}
\end{omfigure}

\section{A link fails}\label{sec:g12:routing:failure}

The program below cuts the link D -- F, then lets both methods find new routes. The distance vectors start from
their stable tables of before; the link-state routers flood the change and run Dijkstra again.

\omcode{code/grade-12/11-routing/listings/failure.py}{1}{17}{The link D -- F fails.}\label{lst:g12:routing:failure}

\omoutput{code/grade-12/11-routing/outputs/failure.txt}{Both methods agree in the end, on a route of cost
$220$ through E. The distance vectors need $13$ rounds, climbing $20$ at a time.}\label{out:g12:routing:failure}

The slow climb is \omterm{def:g12:routing:convergence}{counting to infinity} again, stopped this time by a real route. Right after the cut, D hears E say
``F at $20$'', a route through D itself, and takes it: $30$, through E. E then hears $30$ from D, its \omterm{def:g11:network-protocols:routing-table}{next hop}, and
goes to $40$. For nine rounds D and E each name the other as the \omterm{def:g11:network-protocols:routing-table}{next hop} towards F, and a packet for F would bounce
between them. Only when E's figure passes $110$, the cost of its true route through G, does E switch to it; the news
then needs a few more rounds to reach A.

The link-state routers do not count. Each one learns that the link is gone, removes it from its map, and computes the
new tree in one go. The price is that every router stores the whole map and runs \omterm{met:g12:graph-algorithms:dijkstra}{Dijkstra's algorithm}, where a
distance-vector router stores one line per destination and does additions. The hop-count protocol, with its ceiling
of $15$ hops, was meant only for small networks.
% ledger: g12:rip-2453-diameter

\section{Exercises}\label{sec:g12:routing:exercises}

\begin{exercise}[$\star$]\label{exo:g12:routing:1}
Starting from tables that know only their own router, give D's table after one round of distance vectors on the
network of the chapter, with costs from rates.
\end{exercise}

\begin{exercise}[$\star$]\label{exo:g12:routing:2}
Find the route from A to G that has the fewest hops.
\end{exercise}

\begin{exercise}[$\star$]\label{exo:g12:routing:3}
Find the route from A to G of smallest cost, and its cost.
\end{exercise}

\begin{exercise}[$\star$]\label{exo:g12:routing:4}
Using E's table, and then the tables of the routers on the way, give the route of a packet from E to A.
\end{exercise}

\begin{exercise}[$\star$]\label{exo:g12:routing:5}
Write \code{dv_update(table, neighbour, neighbour_table, cost)}, which returns a router's table after it has heard
one neighbour. Tables map each destination to a pair (distance, \omterm{def:g11:network-protocols:routing-table}{next hop}).
\end{exercise}

\begin{exercise}[$\star\star$]\label{exo:g12:routing:6}
In \code{dv_round}, the comparison \code{<} is changed to \code{>}. What happens on the network of the chapter?
\end{exercise}

\begin{exercise}[$\star\star$]\label{exo:g12:routing:7}
How many rounds does a line of $n$ routers need, starting from tables that know nothing, before its tables stop
changing?
\end{exercise}

\begin{exercise}[$\star\star$]\label{exo:g12:routing:8}
Write \code{cost_from_rate(rate, ref)}. What should it return for a link faster than the reference, and why?
\end{exercise}

\begin{exercise}[$\star\star$]\label{exo:g12:routing:9}
The link B -- D fails instead of D -- F. What is the new route from A to G, and its cost?
\end{exercise}

\begin{exercise}[$\star\star$]\label{exo:g12:routing:10}
A cable can fail in one direction only: D still receives from E, but E no longer receives from D. D keeps listing the
link D -- E; E removes it. A link-state router builds its map from every list it receives, keeping a link when either
end lists it. What goes wrong, and what is the fix?
\end{exercise}

\begin{exercise}[$\star\star$]\label{exo:g12:routing:11}
How many messages does flooding all the link lists send on a ring of $n$ routers?
\end{exercise}

\begin{exercise}[$\star\star\star$]\label{exo:g12:routing:12}
Prove that after $k$ rounds of \code{dv_round}, starting from tables that know only their own router, each router's
distance to each destination is the smallest cost of a route of at most $k$ links, and that destinations with no
such route are absent.
\end{exercise}

\begin{exercise}[$\star\star\star$]\label{exo:g12:routing:13}
In the count of the chapter, A starts at $2$ and B at $1$. Show that both reach the ceiling $M$ after $M - 1$ rounds.
How many rounds would a ceiling of $64$ take? If each round took the $30$ seconds between two periodic updates of the
hop-count protocol, how long would the count of the chapter last?
% ledger: g12:rip-2453-update
\end{exercise}

\begin{exercise}[$\star\star\star$]\label{exo:g12:routing:14}
Change \code{dv_round} to use poisoned reverse, and count the rounds after the failure on the line of three routers,
and on a line of four, R1 -- R2 -- R3 -- R4, when R3 -- R4 fails.
\end{exercise}

\begin{exercise}[$\star\star\star$]\label{exo:g12:routing:15}
Explain why the ceiling of $16$ limits the networks that the hop-count protocol can serve to routes of at most $15$
hops, and why its designers did not choose a much larger ceiling.
\end{exercise}

\section{Problem: Counting to Infinity}\label{sec:g12:routing:problem}

\begin{problem}[{Weekend problem --- how long does bad news take?}]\label{pb:g12:routing:1}
This problem follows the line A -- B -- C of section \ref{sec:g12:routing:infinity}, with \omterm{def:g12:routing:metric}{hop counts}, through a
failure, with and without split horizon, then with link state. The outputs
\ref{out:g12:routing:infinity} and \ref{out:g12:routing:split} give the rounds.

\medskip\noindent\textbf{Part I --- The line.}
\begin{enumerate}
  \item Write the table of each router (destination, distance, \omterm{def:g11:network-protocols:routing-table}{next hop}) once the protocol has \omterm{def:g12:routing:convergence}{converged}.
  \item How many rounds does it take to converge from tables that know nothing?
  \item Why does A need B to know anything about C?
  \item In round $1$ after the failure, B hears ``C at $2$'' from A. What does that $2$ really describe?
\end{enumerate}

\medskip\noindent\textbf{Part II --- Without split horizon.}
\begin{enumerate}[resume]
  \item Explain the first round after the failure: why does B go to $3$, while A stays at $2$?
  \item Explain the second round: why does A go to $4$, although $4$ is larger than its old $2$?
  \item During the count, where does a packet from A to C go?
  \item Why do the values climb by two every two rounds?
  \item At which round do both reach $16$? Check it against the output.
  \item With updates every $30$ seconds, how long does the count last?
\end{enumerate}

\medskip\noindent\textbf{Part III --- With split horizon.}
\begin{enumerate}[resume]
  \item Which line of its table does A leave out when it writes to B? Why?
  \item Why does B mark C unreachable in round $1$?
  \item Why does A need a second round?
  \item Add a router D linked to both A and B, giving a triangle A -- B -- D with C hanging from B. Describe how B
    can still learn a false route to C through D after the failure.
  \item Does poisoned reverse change the answer of the line? Of the triangle?
\end{enumerate}

\medskip\noindent\textbf{Part IV --- With link state.}
\begin{enumerate}[resume]
  \item After the failure, which routers send a new list, and what does it say?
  \item How many messages does the flooding of those lists take on the line?
  \item Why can a link-state router not count to infinity?
  \item What does a link-state router pay for that, in memory and in work?
  \item State the result: how many rounds does each method need before A and B know that C is unreachable?
\end{enumerate}
\end{problem}
